% Chapter 3 — System Analysis and Requirements

\chapter{System Analysis and Requirements}
\label{Chapter3}

This chapter establishes the engineering and clinical requirements for Kent-AI. It formalises the functional and non-functional requirements, presents a feasibility study across technical, operational, and clinical dimensions, specifies the hardware and software constraints, and provides a structured justification for every technology selected in the system stack.

%----------------------------------------------------------------------------------------

\section{Requirements Gathering and Analysis}

Classical homeopathic repertorization is fundamentally a semantic matching problem constrained by clinical safety. Requirements were formulated through analysis of Dr. J. T. Kent's prescribing methodology \cite{kent1897}, consultation with homeopathic practitioners, and systematic review of clinical case-taking protocols.

The analysis revealed that practitioners spend 60--75\% of total consultation time mapping colloquial symptoms into taxonomic rubrics. Furthermore, existing digital tools require manual rubric selection, creating severe cognitive friction. Kent-AI was therefore conceived not as an autonomous diagnostic agent, but as an \textbf{AI-augmented clinical decision-support workstation} that accelerates rubric discovery while keeping the physician in absolute control of final remedy selection.

%----------------------------------------------------------------------------------------

\section{Functional Requirements}

The functional requirements (FR) define the observable behaviours and capabilities of the system:

\begin{enumerate}
  \item \textbf{FR1: Conversational Clinical Intake}
  \begin{itemize}
    \item The system shall guide the patient through a multi-turn natural language dialogue covering seven Kentian dimensions: Mental/Emotional, Anatomical Location, Sensation, Aggravation Modalities, Amelioration Modalities, Concomitants, and Temporal Patterns.
    \item The system shall provide context-aware dynamic quick-reply suggestion chips streamed from an LLM intent classifier.
    \item The system shall support free-form typing, guided multi-choice replies, and paste-in of pre-written clinical consultation notes.
  \end{itemize}

  \item \textbf{FR2: Seven-Dimension Named Entity Recognition (HomeoNER)}
  \begin{itemize}
    \item The system shall extract clinical symptom spans from raw dialogue transcripts and classify them into 15 BIO token labels.
    \item The system shall operate with zero character offset drift between raw input text and extracted spans.
  \end{itemize}

  \item \textbf{FR3: Negation and Coreference Resolution}
  \begin{itemize}
    \item The system shall identify explicitly negated symptoms (e.g., \textit{``no nausea''}, \textit{``denies headache''}) and exclude them from retrieval queries.
    \item The system shall resolve anaphoric references (e.g., \textit{``it gets worse in the evening''} resolved to the previously mentioned pain).
  \end{itemize}

  \item \textbf{FR4: Sub-Second Dense Semantic Rubric Retrieval}
  \begin{itemize}
    \item The system shall index all 74,513 rubrics from Kent's Repertory in a persistent vector database.
    \item Given an extracted symptom entity, the system shall retrieve the Top-$K$ ($K=20$) semantically similar rubrics in under 20 milliseconds.
    \item The system shall support repertory section filtering (e.g., scoping search to MIND rubrics).
  \end{itemize}

  \item \textbf{FR5: Deterministic Grade-Weighted Totality Ranking}
  \begin{itemize}
    \item The system shall rank candidate remedies using Kent's bibliographic grades: Grade 3 (Bold) = 3.0, Grade 2 (Italic) = 2.0, Grade 1 (Roman) = 1.0.
    \item Scores shall incorporate cosine similarity from dense retrieval and inverse rubric frequency specificity.
    \item Dimension weights shall prioritize Mental symptoms (3$\times$) over Modalities (2$\times$), Sensations (2$\times$), and Locations (1$\times$).
  \end{itemize}

  \item \textbf{FR6: Physician Review and Case Editing}
  \begin{itemize}
    \item The clinician shall be able to inspect, edit, add, or delete extracted symptom entities before rubric retrieval.
    \item The clinician shall be able to toggle rubrics between active, muted, and eliminatory (hard filter) modes in an interactive Totality Matrix.
  \end{itemize}

  \item \textbf{FR7: Multi-Page Clinical PDF Export}
  \begin{itemize}
    \item The system shall export a formatted, multi-page clinical report including patient vitals, 7-dimension totality grids, candidate remedy scores, totality score bar charts, and running ``Page X of Y'' headers/footers.
  \end{itemize}

  \item \textbf{FR8: Persistent Case Library Management}
  \begin{itemize}
    \item The system shall persist consultation records, transcripts, totality selections, and final reports in a local SQLite database (\texttt{cases.sqlite}).
    \item The system shall support case filtering, search, pagination, duplicate, and batch JSON/CSV export.
  \end{itemize}
\end{enumerate}

%----------------------------------------------------------------------------------------

\section{Non-Functional Requirements}

\begin{enumerate}
  \item \textbf{NFR1: Retrieval Latency}
  \begin{itemize}
    \item 99th percentile (p99) retrieval latency for a single symptom query against the 74,513-rubric index shall not exceed 25 milliseconds on consumer hardware.
  \end{itemize}

  \item \textbf{NFR2: Clinical Data Privacy and Air-Gapped Operation}
  \begin{itemize}
    \item All patient data, transcripts, embeddings, and LLM inferences shall execute 100\% on-premises without external API calls or cloud egress.
    \item The system shall operate fully offline once local model weights and the vector index are initialized.
  \end{itemize}

  \item \textbf{NFR3: Extraction Accuracy and Reliability}
  \begin{itemize}
    \item LLM clinical extraction shall maintain $\geq$95\% JSON schema validity and $\geq$90\% verbatim clinical grounding on patient narratives.
    \item The system shall never hallucinate nonexistent remedies or alter Kent's published bibliographic grades.
  \end{itemize}

  \item \textbf{NFR4: Usability and Ergonomics}
  \begin{itemize}
    \item The interface shall adhere to modern design tokens: dark slate palette, clear contrast ratios (WCAG AA compliant), 2px focus-visible rings for keyboard accessibility, and zero emojis in clinical telemetry.
    \item The workstation shall enforce a predictable 4-step stepper (Intake $\to$ Review $\to$ Analysis $\to$ Report).
  \end{itemize}

  \item \textbf{NFR5: Maintainability and Test Coverage}
  \begin{itemize}
    \item The codebase shall maintain $\geq$90\% test pass rate with fully isolated unit tests using mock backends, enabling test suite execution without GPU or Ollama dependencies in under 30 seconds.
  \end{itemize}
\end{enumerate}

%----------------------------------------------------------------------------------------

\section{Feasibility Study}

\subsection{Technical Feasibility}
The system relies on transformer-based sequence labeling, dense vector indexing, and quantized LLM inference. High-performance biomedical transformers (PubMedBERT, BioLinkBERT) are freely available via HuggingFace. ChromaDB provides an embedded HNSW implementation requiring no external server process. Local LLM execution via Ollama (LLaMA 3 8B, BioMistral 7B) requires 8--16 GB VRAM, which is readily available on modern workstations and institutional GPU clusters. Technical feasibility is confirmed by our working implementation.

\subsection{Operational Feasibility}
The workstation conforms to the standard clinical intake workflow of homeopathic practitioners: case-taking, symptom recording, totality analysis, and repertorization. By automating the transcription-to-rubric mapping while retaining the practitioner's ultimate decision-making authority, the system integrates smoothly into existing clinical routines without imposing artificial constraints.

\subsection{Economic and Licensing Feasibility}
All system components --- Python, PyTorch, HuggingFace Transformers, ChromaDB, Ollama, Streamlit, SQLite, and ReportLab --- are released under permissive open-source licenses (MIT, Apache 2.0, BSD). The Kent's Repertory dataset is derived from the public-domain 1897 edition. The system incurs zero ongoing per-query API costs.

\subsection{Clinical and Regulatory Feasibility}
Kent-AI is designated as an \textbf{Investigational Clinical Decision Support Tool} under Section 520(o) of the US FDA 21st Century Cures Act guidelines for non-device CDS software. It displays clear disclaimers, surfaces source-grounded rubric provenance, and does not provide automated diagnoses or prescriptions.

%----------------------------------------------------------------------------------------

\section{Hardware and Software Requirements}

Table \ref{tab:hardwaresoftware} outlines the hardware and software constraints for development, training, and production deployment.

\begin{table}[h]
\caption[System Hardware and Software Constraints]{System Hardware and Software Specifications.}
\label{tab:hardwaresoftware}
\centering
\small
\begin{tabular}{lll}
\toprule
\tabhead{Environment} & \tabhead{Component} & \tabhead{Specification} \\
\midrule
\multirow{4}{*}{\textbf{Training / GPU Server}} & GPU & NVIDIA RTX 5000 Ada (32GB GDDR6 ECC) \\
 & CUDA / Driver & CUDA 12.4, NVIDIA Driver $\geq$ 550.0 \\
 & CPU / RAM & 64-core Intel Xeon, 128 GB DDR5 ECC \\
 & Storage & 2 TB NVMe SSD \\
\midrule
\multirow{4}{*}{\textbf{Inference / Deployment}} & GPU (Optional) & NVIDIA RTX 3060 / 4060 (8--16 GB VRAM) \\
 & CPU Mode & Intel Core i7 / AMD Ryzen 7 (16+ GB RAM) \\
 & Storage & 15 GB available disk space (models + index) \\
 & OS & Linux (Ubuntu 22.04+), Windows 11, macOS (Metal) \\
\midrule
\multirow{4}{*}{\textbf{Software Stack}} & Language & Python 3.12 (standard virtual environment) \\
 & Deep Learning & PyTorch 2.6.0, Transformers 4.46, PEFT 0.14 \\
 & Vector Store & ChromaDB 0.5.5, sentence-transformers 3.0 \\
 & UI / PDF & Streamlit 1.64.0, ReportLab 5.1, PyMuPDF 1.24 \\
\bottomrule
\end{tabular}
\end{table}

%----------------------------------------------------------------------------------------

\section{Tools and Technologies Chosen with Justification}

Table \ref{tab:techjustification} details the engineering justification for every major technology selected, comparing each against rejected alternatives.

\begin{table}[h]
\caption[Technology Selection Justification]{Detailed Justification for Technology Stack Choices.}
\label{tab:techjustification}
\centering
\scriptsize
\begin{tabular}{p{2.5cm}p{3.2cm}p{3.5cm}p{4cm}}
\toprule
\tabhead{Component} & \tabhead{Selected Technology} & \tabhead{Evaluated Alternatives} & \tabhead{Justification} \\
\midrule
\textbf{Vector Store} & ChromaDB (HNSW) & FAISS, Milvus, Qdrant, pgvector & Embedded Python native, zero daemon setup, built-in persistence, native metadata filtering on section IDs, sub-18ms latency. \\
\midrule
\textbf{Sentence Embedder} & \texttt{all-MiniLM-L6-v2} & OpenAI \texttt{text-embedding-ada-002}, Clinical-BERT & 384-dim compact unit-normalized vectors, fast CPU/GPU inference ($<$5ms), zero cloud egress, completely open-source. \\
\midrule
\textbf{Biomedical Encoders} & BioLinkBERT \& PubMedBERT & Standard BERT, RoBERTa, ClinicalBERT alone & PubMedBERT trained from scratch on biomedical texts (no subword fragmentation); BioLinkBERT captures citation-graph modality linkages. \\
\midrule
\textbf{LLM Fine-Tuning} & QLoRA (NF4 + BitsAndBytes) & Full fine-tuning, Prefix tuning, LoRA-16bit & Reduces 7B model VRAM from 28GB to $\sim$6GB, enabling fine-tuning on a single workstation GPU with minimal performance loss. \\
\midrule
\textbf{Local LLM Runtime} & Ollama REST API & vLLM, TGI, llama.cpp direct bindings & Turnkey model management, robust streaming HTTP REST API, standardized JSON schema enforcement, cross-platform. \\
\midrule
\textbf{Frontend Workstation} & Streamlit ($\geq$1.64.0) & React + FastAPI, Electron, PyQt & Fast clinical iteration, built-in session state management, native Python integration with ML models, fragment reruns for snappy UI. \\
\midrule
\textbf{Clinical PDF Engine} & ReportLab 5.1 + PyMuPDF & Weasyprint, wkhtmltopdf, LaTeX CLI & Deterministic canvas rendering, exact two-pass page numbering (\texttt{NumberedCanvas}), native vector tables and charts, zero external binary dependency. \\
\midrule
\textbf{Case Database} & SQLite 3 & PostgreSQL, MongoDB & Embedded zero-admin ACID storage, file-level portability (\texttt{cases.sqlite}), seamless integration with Python's standard library. \\
\bottomrule
\end{tabular}
\end{table}

